\begin{table}[htbp]
  \caption{Comparison of spatial models for uncertain-onset stroke}
  \label{tab:spatial_models}
  \small
  \begin{threeparttable}

    \begin{tabularx}{\textwidth}{@{}Xcccc@{}}
      \toprule
      \textbf{Model}      &
      \textbf{Spatial SD} &
      \textbf{95\% CrI}   &
      \boldmath$\phi$     &
      \textbf{DIC / WAIC}                                                      \\
      \midrule

      Unadjusted          & 0.696 & 0.614--0.781 & 0.976 & 61088.99 / 61089.08 \\

      Sociodemographic    & 0.697 & 0.616--0.783 & 0.976 & 60753.25 / 60752.12 \\

      Clinical            & 0.700 & 0.618--0.786 & 0.975 & 60780.37 / 60782.12 \\

      Missingness         & 0.775 & 0.683--0.871 & 0.969 & 54270.76 / 54269.53 \\

      Health trust        & 0.518 & 0.444--0.595 & 0.970 & 60710.35 / 60710.54 \\

      Full model          & 0.612 & 0.529--0.670 & 0.969 & 53638.95 / 53640.24 \\

      \bottomrule
    \end{tabularx}

    \begin{tablenotes}
      \footnotesize
      \item Spatial SD - the posterior standard deviation of the spatial effect. DIC =
      deviance information criterion; WAIC = Watanabe-Akaike information criterion.
      $\phi$ - the proportion of spatial variance attributed to the structured
      component.
    \end{tablenotes}

  \end{threeparttable}
\end{table}